\documentclass[12pt]{article}
\usepackage{amsmath, amssymb, amsthm}
\usepackage[english]{babel}
\usepackage{geometry}
\usepackage{mathrsfs}

\numberwithin{equation}{section}

% Other packages...
\newtheorem{theorem}{Theorem}[section]
\newtheorem{corollary}{Corollary}[theorem]
\newtheorem{lemma}[theorem]{Lemma}

\title{\textbf{Exercise Solutions of Chp 2 N. Goldenfeld Phase T. and RN}}
\author{Dr. Jed Pixley and Elisha Shmalo}
\date{\today}

\geometry{left=1in, right=1in, top=1in, bottom=1in}
\renewcommand{\baselinestretch}{1.25}

\begin{document}

\maketitle

\section{Chapter 2}

\textbf{Exercise 2.1} Concerning the free energy of the n-n Ising model:

\begin{equation}
    -H_\Omega (\{S\}) = H \sum_i S_i + J\sum_{i,j}S_iS_j.
\end{equation}

(a) We show that $g(\beta) = \lim_{N(\Omega) \to \infty}\log(Z_\Omega(\beta))/N(\Omega)$ is convex down in $\beta$.

Recall 

\begin{equation}
    Z(\beta) = \sum_{\{S_i\}}e^{-\beta H_\Omega (\{S\})} = \sum_{\{S_i\}}e^{\beta H \sum_i S_i + \beta J\sum_{i,j}S_iS_j}.
\end{equation}

Proceeding dead ahead, with the definition of convex down in mind, let $\beta_2 > \beta_1$ and $a,b \in [0, 1]$ such that $a + b = 1$ and consider 

\begin{equation}
    Z_\Omega(a\beta_1 + b\beta_2) = \sum_{\{S_i\}} [e^{(\beta_1 H \sum_i S_i + \beta_1 J\sum_{i,j}S_iS_j)^a}e^{(\beta_2 H \sum_i S_i + \beta_2 J\sum_{i,j}S_iS_j)^b}]
\end{equation}

Now by Holder's inequality this can be bounded above as 

\begin{equation}
    \begin{aligned}
        \sum_{\{S_i\}} [e^{(\beta_1 H \sum_i S_i + \beta_1 J\sum_{i,j}S_iS_j)^a}e^{(\beta_2 H \sum_i S_i + \beta_2 J\sum_{i,j}S_iS_j)^b}] &\leq (\sum_{\{S_i\}} e^{(\beta_1 H \sum_i S_i + \beta_1 J\sum_{i,j}S_iS_j)})^a \\ 
        & \times (\sum_{\{S_i\}} e^{(\beta_2 H \sum_i S_i + \beta_2 J\sum_{i,j}S_iS_j)})^b
    \end{aligned}
\end{equation}

But this is precisely the same as 

\begin{equation}
    Z_\Omega(a\beta_1 + b\beta_2) \leq Z_\Omega(\beta_1)^aZ_\Omega(\beta_2)^b.
\end{equation}

It follows that 

\begin{equation}
    \log(Z_\Omega(a\beta_1 + b\beta_2)) \leq a\log(Z_\Omega(\beta_1)) + b\log(Z_\Omega(\beta_2))
\end{equation}

We then see imediatly that 

\[
\begin{aligned}
    g(a\beta_1 + b\beta_2) &= \lim_{N(\Omega) \to \infty}\log(Z_\Omega(a\beta_1 + b\beta_2))/N(\Omega)  \\
    & \leq \lim_{N(\Omega) \to \infty}(a\log(Z_\Omega(a\beta_1))/N(\Omega) + b\log(Z_\Omega(a\beta_1))/N(\Omega)) \\
    &= ag(\beta_1) + bg(\beta_2).
\end{aligned}
\]

and we are done.

(b) We will show hat
the free energy per unit volume

\begin{equation}
    f(T) = -\frac{1}{\beta}g(\beta) = -k T g(1/(kT))
\end{equation}

is concave up in T.

\textbf{Ask Professor}

(c) We have the Gibbs free energy defined as 

\begin{equation}
    \Gamma_\Omega(M) = F_\Omega(H(M)) + N(\Omega)MH(M)
\end{equation}

where $H(M)$ is defined implicitly as by $M(H) = -\partial f / \partial H$. We will show that $\Gamma(M)$ is 
convex down in $M$.

\end{document}